\chapter{总结与展望}\label{ch:conclusion}

\section{本文工作总结}

随着5G网络的全面商用以及6G技术的快速发展，人们对网络的覆盖能力提出了更高的要求。传统的适应信道的设计范式已难以满足全域覆盖的要求，可重构天线系统可以主动调控电磁传播环境，实现信道重构，是6G物理层的一个重要发展方向。
本论文围绕动态超表面天线、捏合天线系统和可重构智能表面三类典型的可重构天线技术，针对DMA存在的幅相耦合约束和离散相移受限的特征、PASS通信与感知的构型冲突以及RIS盲波束赋形开销大且缺乏实测等挑战，从闭式性能分析、波束赋形算法设计、原型系统研制三个层面完善了可重构天线系统相关理论，论文的主要研究内容和结论归纳如下：

针对DMA系统中洛伦兹幅相耦合导致松弛投影方法存在相位失配、考虑波导衰减下的性能分析缺失的问题，本文研究了连续相移DMA的波束赋形性能分析与相位对齐算法。首先，通过解耦洛伦兹约束的固定项和可调项，推导出了单用户下的最优相位闭式解。在此基础上，分析了系统的平均信噪比、中断概率、遍历容量等性能指标，得出了系统的波束赋形增益随着单波导单元数量呈现先增后降的结论，给出了最优单元数量的近似解。随后，定量对比了不同投影策略的归一化增益，从理论上证明了解耦求解的方法优于松弛方法，还分析了CSI不完美造成的性能损失。为降低大规模DMA波束赋形算法对CSI的依赖，提出了两种基于RSSI反馈的波束赋形算法。其中，RECD算法通过逐元素坐标下降与线性回归求解，逐步迭代逼近最优解；NI-REPA算法通过引入全局参考信号，剔除了测量过程中的待测单元的干扰，仅需一轮测量即可完成相位对齐。通过对线性回归模型进行误差分析，证明了上述算法在均匀采样下能取得各向同性且最小的均方误差。仿真结果验证了理论推导的准确性，在高发射功率的条件下，NI-REPA算法的归一化接收功率可达到$95\%$以上。


针对实际的DMA系统相移受限于低精度离散控制和极化率约束，呈现出相位分布不均匀、移相能力受限的特征，且直接使用连续解投影会造成性能退化的问题，本文研究了任意离散相移下的性能分析和角域全局寻优算法。首先，在考虑莱斯信道衰落和任意离散相位分布的系统模型下，给出了采用就近投影策略时的系统平均信噪比和量化损失，并推导了随机相移分布下的平均性能作为对比基准。随后，从理论上证明了当移相能力充足时，均匀分布能够最小化量化损失；当移相能力受限时，则应该在尽可能利用动态范围的前提下实现局部均匀分布，同时揭示了1-bit量化时最近投影方法可能失效。在算法方面，为减少CPP算法的性能损失，提出了一种多项式时间复杂度的ADSPS算法。该算法通过引入辅助变量与几何投影，将多变量的离散组合优化问题转化为一维角域上的遍历搜索，并结合临界角矩阵的低秩结构，设计了基于多路归并的分组排序策略，在保证取得全局最优解的前提下，搜索的复杂度相比穷举法大大降低。理论下界表明，在1-bit量化下，该算法能保证至少获得连续相移条件下四分之一的阵列增益。最后，通过不同相移配置的数值仿真验证了ADSPS算法的性能优势。结果表明，对于1-bit均匀量化的DMA阵列，ADSPS算法较CPP算法降低了0.3 dB的量化损失。


针对PASS在大尺度空间重构下通信最优与感知最优在物理几何构型上存在矛盾、缺乏最优定位几何构型理论指导的问题，本文开展了PASS波束赋形性能分析与几何辅助定位算法研究。首先，推导了天线对齐策略下平均信噪比和空间中断概率的表达式，定量刻画了PASS大尺度重构相较固定阵列所具备的空间覆盖优势。随后，在感知定位层面，基于费雪信息矩阵的迹最大化准则，证明了最优的定位天线构型是一个以用户正上方为圆心、半径近似为波导高度的圆。以此圆形构型为基础，结合二维MUSIC空间谱估计，提出了CG-APPL算法，交替执行这两个步骤，逐渐逼近最优的观测几何。仿真结果表明，CG-APPL算法克服了固定阵列在近端径向分辨率缺失和远端几何精度衰减的问题，在整个服务区域内保持了稳定的亚毫米级定位误差，实现了通信与感知性能的提升。


针对现有RIS盲波束赋形算法在NLoS环境下容易失效、随机采样开销过大且寻优的过程阻断数据传输的问题，本文研究了支持在线传输的低复杂度波束赋形算法，研制了工作在Sub-6 GHz频段的大规模RIS原型系统并开展了多场景的现场测试。首先，利用均匀平面阵列的波束结构特性，受导向矢量和DFT矩阵的启发，提出了GFB算法，通过两阶段的贪婪波束赋形，降低了搜索的复杂度。进一步地，为了增强RIS系统在多径环境下的鲁棒性，提出了变步长的ASG算法。在迭代初期，算法会同时改变较多单元的相位，有利于跳出局部最优；迭代后期转入精细化搜索，兼顾了寻优速度与在线传输性能。仿真结果表明，所提ASG算法收敛较快，在不同量化比特数、不同莱斯因子下均具有较高的性能。
在系统实现层面，研制了两套分别工作于5.8 GHz和2.6 GHz的RIS面板，实现了射频收发、基带信号处理、RIS面板控制和RSSI反馈的闭环调试。随后，在实验室环境下验证了RIS的信道互易性，评估了系统的波束调控性能。
现场实测结果表明，在室内穿墙NLoS场景、L型走廊盲区、办公区边缘覆盖场景下，RIS分别实现了最高 $26 \text{ dB}$、$20 \text{ dB}$ 和 $22 \text{ dB}$ 的接收功率增益。最后，在商用5G现网测试中，部署RIS使得室内终端的RSSI提升了 $4 \text{ dB}$ 至 $5 \text{ dB}$，通过对比实验揭示了商用环境中的宽波束基站与全向天线对RIS增益的稀释效应，为后续商用部署提供了实验参考。


\section{下一步工作展望}

本文围绕动态超表面天线、捏合天线系统和可重构智能表面三类可重构天线技术，在波束赋形算法设计、系统性能分析和原型验证等方面取得了一定进展。然而，受限于博士期间的研究时间和实验条件，目前的工作仍存在一定局限与不足：在理论层面，本文所提算法主要面向单用户场景；在系统层面，三类可重构技术的研究相对独立，缺乏联合部署与协同优化的探索；在实测层面，DMA与PASS的波束赋形方案尚未在硬件平台上进行验证。未来研究可以从以下几个方面展开：

\begin{enumerate}
	\item 探究多用户场景下的波束赋形算法设计。在实际 6G 网络中，基站需同时服务大量用户，多用户间的信道干扰将严重影响频谱效率。在此背景下，如何利用DMA和PASS独特的调控架构有效消除多用户间的同频干扰，是一个极具挑战性的问题。后续研究可考虑引入系统和速率最大化准则，设计适用于多用户干扰环境的低复杂度波束赋形算法。
	
	\item 探索可重构天线技术的混合部署与协同优化。本文分别研究了改进发射机架构的DMA和PASS技术以及改善电磁传播环境的RIS技术。单一的技术能带来一定增益，但在复杂多变的6G全域覆盖与通感一体化场景下仍存在局限。后续研究可针对“DMA+RIS”或“PASS+RIS”这两类混合架构，探索联合波束赋形与位置部署的协同优化问题，有望进一步提升通信系统性能。
	
	\item 进行DMA与PASS系统硬件原型的研制与实测验证。受限于实验条件，本文对 DMA 和 PASS 的研究停留在理论分析与数值仿真阶段，所提出的波束赋形算法未在原型系统上得到检验。后续计划研制基于超材料谐振单元和波导馈电的 DMA 原型阵列，并搭建由高精度电机驱动的PASS实验平台，验证本文提出的波束赋形与感知定位算法。随后，评估实际系统中各类非理想因素对系统性能的影响，例如由加工误差引起的DMA幅相响应偏移，以及因步进电机驱动引起的PA移动速度受限和位置偏差等情形，建立包含这些非理想因素的系统模型，并定量分析它们对系统性能的影响。
\end{enumerate}
